\section{Biconditional Statements}
It is important to make a distinction between $P\Rightarrow Q$ and $Q\Rightarrow P$. We say that when $P\Rightarrow Q$, then $Q\Rightarrow P$ is its \emph{converse}. However, when both statements are true we say that it is a \emph{biconditional}. 
$$
(P\Rightarrow Q)\land(Q\Rightarrow P) := P \Leftrightarrow Q
$$
In terms of the \emph{semantics} or, in other words, the truth table of the biconditional, the following is true: 
\begin{table}[H]
    \centering
    \begin{tabular}{ccc}
        \toprule
        $P$ & $Q$ & $P\Leftrightarrow Q$\\
        \midrule
        T & T & T\\
        
        T & F & F \\
        
        F & T & F \\
        
        F & F & T\\
        \bottomrule
    \end{tabular}
\end{table}
We see that the biconditional is correct \emph{only when} $P=T$ and $Q=T$ or when they both equal $F$.\newline
There are many ways of saying `if and only if' in English, some are as follows: 
\begin{enumerate}
    \item $P$ iff $Q$.
    \item $P$ is a necessary and sufficient condition for $Q$. 
    \item $P$ is equivalent to $Q$.
    \item For $P$ it is necessary and sufficient that $Q$. 
    \item If $P$, then $Q$ and conversely. 
\end{enumerate}

\subsection{Exercises for Section 14}
\begin{exercise}\label{ex:2.14}
    Without changing their meanings, convert each of the following sentences into a sentence having the form: $P\Leftrightarrow Q$. 
    \begin{enumerate}
        \item For a matrix $A$ to be invertible, it is necessary and sufficient that $\det(A)\ne0$. 
        \item If a function has a constant derivative, then it is linear, and conversely. 
        \item If $xy=0$, then $x=0$ or $y=0$, and conversely. 
        \item If $a\in\Q$, then $5a\in\Q$, and if $5a\in\Q$, then $a\in\Q$. 
        \item For an occurence to become an adventure, it is necessary and sufficient for one to recount it. 
    \end{enumerate}
\end{exercise}\hyperref[sol:2.14]{[Go to Solution]}